\section{总结与延伸}

\subsection{核心要点回顾}

本讲围绕"如何高效求解扩散模型的逆过程 ODE"展开，核心贡献如下：

\begin{enumerate}
    \item \textbf{半线性结构的发现}：PF-ODE 的右端项可以分解为线性部分（可精确求解）和非线性部分（需要数值近似），这为高效求解提供了基础
    \item \textbf{指数积分器的应用}：通过积分因子法精确处理线性部分，将数值近似限定在噪声预测的积分上
    \item \textbf{$\lambda$ 空间的参数化}：用 log-SNR 做变量替换，大幅简化积分表达式
    \item \textbf{DDIM 的重新诠释}：DDIM 是 DPM-Solver 的一阶特例，而非朴素欧拉法
    \item \textbf{高阶方法的递归构造}：通过在 $\lambda$ 空间对噪声预测做高阶多项式近似，以少量额外 NFE 换取显著精度提升
    \item \textbf{DPM-Solver++}：转为对数据预测做多项式近似，进一步提升有限步精度
\end{enumerate}

\subsection{更广阔的视角}

\begin{knowledgebox}{为什么数学理解如此重要}
讲者在课程中反复强调：所有这些加速都是在\textbf{不改变训练模型}的前提下实现的。这深刻体现了一个原则——对问题结构的数学理解可以直接转化为实际性能提升。从 DDPM 到 DDIM 到 DPM-Solver，每一步都是对同一个扩散过程的更深入理解。
\end{knowledgebox}

\subsection{后续发展}

DPM-Solver 之后，扩散模型的加速采样继续快速发展：

\begin{itemize}
    \item \textbf{一致性模型}（Consistency Models, Song et al., 2023）：学习直接从噪声跳到数据的映射，可实现单步生成
    \item \textbf{渐进蒸馏}（Progressive Distillation, Salimans \& Ho, 2022）：通过知识蒸馏将多步过程压缩为少步
    \item \textbf{Flow Matching}（Lipman et al., 2022）：设计更直的传输路径，使得简单的欧拉法即可高效工作
    \item \textbf{Rectified Flow}（Liu et al., 2022）：通过 Reflow 操作使 ODE 轨迹趋于直线
\end{itemize}

\subsection{拓展阅读}

\begin{itemize}
    \item Lu, C., et al. (2022). \textit{DPM-Solver: A Fast ODE Solver for Diffusion Probabilistic Model Sampling in Around 10 Steps}. NeurIPS 2022.
    \item Lu, C., et al. (2022). \textit{DPM-Solver++: Fast Solver for Guided Sampling of Diffusion Probabilistic Models}. arXiv:2211.01095.
    \item Song, J., et al. (2020). \textit{Denoising Diffusion Implicit Models}. ICLR 2021.
    \item Song, Y., et al. (2021). \textit{Score-Based Generative Modeling through Stochastic Differential Equations}. ICLR 2021.
    \item Karras, T., et al. (2022). \textit{Elucidating the Design Space of Diffusion-Based Generative Models}. NeurIPS 2022.
\end{itemize}

%% --- Fin del contenido --- %%

\end{document}
